% DEIXIS evidence report, IEEEtran export for XeLaTeX.
% Compile with: latexmk -xelatex report-synthetic-report-v12.tex
% Needs XeLaTeX, BibTeX and the IEEEtran class and IEEEtran.bst style; they are not included.
% Export notes: 1
% 1. Table I was split into 2 column groups of at most 7 data columns.

\documentclass[journal]{IEEEtran}
\usepackage{fontspec}
\usepackage{amsmath}
\usepackage{amssymb}
\usepackage{bm}
\usepackage{xcolor}
\usepackage{cite}
\usepackage{url}
\usepackage{booktabs}
\usepackage{longtable}
\usepackage{array}

\begin{document}

\renewcommand{\abstractname}{Özet}
\renewcommand{\IEEEkeywordsname}{Dizin Terimleri}
\renewcommand{\refname}{Kaynaklar}

\title{SYNTHETIC report \& \ensuremath{\alpha}}

\maketitle

\noindent Kanıt raporu \ensuremath{\cdot} V12

\section*{I. Giriş}

SYNTHETIC claim. \cite{Synth26}

\section*{IV. Literatür Sentezi}

SYNTHETIC table follows.

\onecolumn
{\scriptsize
\begin{longtable}{>{\raggedright\arraybackslash}p{1.0in}>{\raggedright\arraybackslash}p{\dimexpr(\textwidth-1.0in-2\tabcolsep)/7-2\tabcolsep\relax}>{\raggedright\arraybackslash}p{\dimexpr(\textwidth-1.0in-2\tabcolsep)/7-2\tabcolsep\relax}>{\raggedright\arraybackslash}p{\dimexpr(\textwidth-1.0in-2\tabcolsep)/7-2\tabcolsep\relax}>{\raggedright\arraybackslash}p{\dimexpr(\textwidth-1.0in-2\tabcolsep)/7-2\tabcolsep\relax}>{\raggedright\arraybackslash}p{\dimexpr(\textwidth-1.0in-2\tabcolsep)/7-2\tabcolsep\relax}>{\raggedright\arraybackslash}p{\dimexpr(\textwidth-1.0in-2\tabcolsep)/7-2\tabcolsep\relax}>{\raggedright\arraybackslash}p{\dimexpr(\textwidth-1.0in-2\tabcolsep)/7-2\tabcolsep\relax}}
\caption{Bu rapor için dondurulan kanıt tablosu (2 kaynak, 10 sütun). (10 sütundan 1–7)}\\
\toprule
Kaynak & SYNTHETIC C1 & SYNTHETIC C2 & SYNTHETIC C3 & SYNTHETIC C4 & SYNTHETIC C5 & SYNTHETIC C6 & SYNTHETIC C7 \\
\midrule
\endfirsthead
\caption[]{Bu rapor için dondurulan kanıt tablosu (2 kaynak, 10 sütun). (10 sütundan 1–7) (devamı)}\\
\toprule
Kaynak & SYNTHETIC C1 & SYNTHETIC C2 & SYNTHETIC C3 & SYNTHETIC C4 & SYNTHETIC C5 & SYNTHETIC C6 & SYNTHETIC C7 \\
\midrule
\endhead
\bottomrule
\endlastfoot
{}[1] Synth26 & 2 m & 3 & Evet & Hayır & SYNTHETIC option, unknown &  & unverified (doğrulanmadı: bağlı alıntı yok) \\
SYNTHETIC row 2 & SYNTHETIC R2C1 & SYNTHETIC R2C2 & SYNTHETIC R2C3 & SYNTHETIC R2C4 & SYNTHETIC R2C5 & SYNTHETIC R2C6 & SYNTHETIC R2C7 \\
\end{longtable}}
\addtocounter{table}{-1}
{\scriptsize
\begin{longtable}{>{\raggedright\arraybackslash}p{1.0in}>{\raggedright\arraybackslash}p{\dimexpr(\textwidth-1.0in-2\tabcolsep)/3-2\tabcolsep\relax}>{\raggedright\arraybackslash}p{\dimexpr(\textwidth-1.0in-2\tabcolsep)/3-2\tabcolsep\relax}>{\raggedright\arraybackslash}p{\dimexpr(\textwidth-1.0in-2\tabcolsep)/3-2\tabcolsep\relax}}
\caption{Bu rapor için dondurulan kanıt tablosu (2 kaynak, 10 sütun). (10 sütundan 8–10)}\\
\toprule
Kaynak & SYNTHETIC C8 & SYNTHETIC C9 & SYNTHETIC C10 \\
\midrule
\endfirsthead
\caption[]{Bu rapor için dondurulan kanıt tablosu (2 kaynak, 10 sütun). (10 sütundan 8–10) (devamı)}\\
\toprule
Kaynak & SYNTHETIC C8 & SYNTHETIC C9 & SYNTHETIC C10 \\
\midrule
\endhead
\bottomrule
\endlastfoot
{}[1] Synth26 & (doğrulanmadı: bağlı alıntı yok) & not found & — \\
SYNTHETIC row 2 & SYNTHETIC R2C8 & SYNTHETIC R2C9 & SYNTHETIC R2C10 \\
\end{longtable}}
\twocolumn

\section*{Rapor notları}

DEIXIS ile üretildi.

Atıf çapaları ilgili pasajlarda veya hücrelerde bulundu. Bu rapor için model veya kişi incelemesi kaydedilmedi; kod, her pasajın ilgili iddiayı destekleyip desteklemediğini denetlemedi.

Bu çıktı raporun metnini ve kaynakçasını taşır. Alıntı pasajları, tablo hücresi kanıtları ve konumlanmış atıf çapaları DEIXIS'te kalır.

\bibliographystyle{IEEEtran}

\bibliography{report-synthetic-report-v12}

\end{document}
